\chapter{Background and Related Work}
\label{ch:background}

\section{Spiking neural networks}

A spiking neural network communicates through discrete binary events rather than
continuous activations. The neuron model used throughout this work is the leaky
integrate-and-fire (LIF) unit in its discrete-time form: the membrane potential
decays by a factor $\beta$ at each step, integrates the incoming current, emits a
spike when it crosses a threshold, and is reset.

\begin{equation}
U[t+1] = \beta U[t] + W X[t+1] - S[t] U_{\text{thr}},
\qquad
S[t] = \begin{cases}
1 & \text{if } U[t] > U_{\text{thr}} \\
0 & \text{otherwise}
\end{cases}
\label{eq:lif}
\end{equation}

Two properties follow from Equation~\ref{eq:lif} and shape everything in this
thesis. First, the activation $S[t]$ is binary, so a synaptic operation is an
accumulation rather than a multiply-accumulate; this is the origin of the energy
advantage discussed in Section~\ref{sec:energy-background}. Second, $U[t]$ is a
\emph{state}: it must be stored between time steps for every neuron in the
network, which makes it the dominant memory cost at deployment and a
quantization target in its own right.

\subsection{Training by surrogate gradients}

The spike-generating function in Equation~\ref{eq:lif} is a Heaviside step, whose
derivative is zero almost everywhere and undefined at threshold. Backpropagation
through time is therefore not directly applicable.

The standard remedy is the \emph{surrogate gradient}~\cite{neftci2019surrogate}:
the forward pass keeps the exact discontinuous spike, while the backward pass
substitutes the derivative of a smooth approximant. The choice of approximant
--- sigmoid, fast sigmoid, arctangent --- affects convergence but not the
forward computation, and networks trained this way reach accuracies comparable
to conventional architectures on image benchmarks~\cite{eshraghian2023}.

The substitution has a consequence for explainability that motivates the choice
of method in Section~\ref{sec:sam-background}: any gradient-based attribution
computed on such a network inherits the approximation error of the surrogate,
because the gradients it uses are not the gradients of the function the network
actually computes.

\subsection{Input encoding}

Spiking networks require an input that varies over time, while images do not. The
two common strategies are \emph{rate coding}, in which pixel intensities become
Bernoulli firing probabilities, and \emph{direct} or \emph{constant} encoding, in
which the analog frame is presented identically at every time step and the first
convolutional layer acts as a learned analog-to-spike converter.

Direct encoding, introduced in this form by DIET-SNN~\cite{rathi2023dietsnn},
converges in fewer time steps because it does not inject sampling noise into the
input. It carries an interpretive consequence that must be stated explicitly:
since the stimulus does not vary, \textbf{all temporal structure in the network
originates from the internal LIF dynamics} --- leak, integration and reset ---
rather than from variation in the input. Every temporal analysis in this thesis
is to be read in that light.

\section{Quantization}

Quantization maps a continuous parameter to one of $2^b$ discrete levels. For a
symmetric uniform quantizer over $[-w_{\max}, +w_{\max}]$ the step is
\begin{equation}
\Delta = \frac{2 w_{\max}}{2^b - 1},
\label{eq:step}
\end{equation}
and the maximum rounding error is $\Delta/2$. Equation~\ref{eq:step} is used in
Section~\ref{sec:perturbation} to express a perturbation in units the
quantization itself defines.

Two families of technique exist. \emph{Post-training quantization} rounds the
weights of a trained model, and is what Rabbi et al.~\cite{rabbi2026} apply.
\emph{Quantization-aware training} simulates the rounding during training through
fake-quantization, and propagates gradients through the non-differentiable
rounding operation with a straight-through estimator~\cite{bengio2013ste} ---
the same device, in a different guise, as the surrogate gradient of
Section~2.1.1.

\subsection{Quantization for spiking networks}

Spiking architectures add constraints of their own. The membrane decay $\beta$
multiplies a stored state at every time step, so on integer hardware it is
cheapest when realizable as a bit shift: $\beta = 2^{-k}$ is a right shift, and
$\beta = 1 - 2^{-k}$ is a shift plus a subtraction. MINT~\cite{yin2024mint}
additionally shares a single quantization scale between weights and membrane
potentials, which eliminates multipliers from the inference loop entirely.

Smith et al.~\cite{smith2026} argue independently that accuracy is an
insufficient criterion for quantized spiking networks, and propose an Earth
Mover's Distance between per-neuron firing-rate distributions as a diagnostic of
behavioural divergence. They separate weight from membrane quantization as this
thesis does, and report that membrane-potential quantization requires at least
four bits to preserve both accuracy and dynamics while learned weight
quantization remains viable at two. They ground the asymmetry mechanistically: a
weight quantization error is fixed for an entire forward pass, whereas a membrane
error propagates through the recurrent update at every time step.

\subsection{Energy accounting}
\label{sec:energy-background}

Energy is estimated analytically rather than measured, as is standard in the
spiking literature. The count of synaptic operations is multiplied by the
per-operation cost at a reference process node, conventionally the 45\,nm
figures of Horowitz~\cite{horowitz2014}: a 32-bit floating-point
multiply-accumulate costs $4.6$\,pJ against $0.9$\,pJ for a plain accumulate.
That factor of five is the physical reason a spiking network can be cheaper,
since a binary activation removes the multiplication, and the accumulation is
performed only where a spike occurred.

\section{Explainability for spiking networks}
\label{sec:sam-background}

Attribution methods developed for convolutional networks transfer badly to
spiking ones. Gradient-based methods such as Grad-CAM~\cite{selvaraju2017gradcam}
require gradients that, as noted above, exist only as surrogates. Perturbation
methods such as LIME are architecture-agnostic but expensive, requiring hundreds
of forward passes per explanation.

The Spike Activation Map (SAM) of Kim and Panda~\cite{kim2021sam} was designed
natively for spiking networks and requires neither gradients nor labels. It is
built on the observation that a neuron firing repeatedly at short intervals is
contributing more to the current decision than one firing sporadically. A
\emph{temporal spike contribution score} weights a past spike at $t'$ by its
recency,
\begin{equation}
T(t, t') = \exp\left(-\gamma\,|t - t'|\right),
\end{equation}
a \emph{neuronal contribution score} sums those weights over a neuron's
preceding spikes, and the map at time $t$ masks that score with the spikes
actually emitted at $t$. The result is one spatial map per time step: an
explanation with a temporal axis, which is what makes the question of this thesis
answerable at all.

Feature-attribution methods for spiking networks have also been surveyed more
broadly~\cite{nguyen2023attribution}, but SAM remains the method that reads the
spike trains directly.

\subsection{Evaluating an explanation}

An explanation that is stable can still be wrong, so a separate faithfulness
test is required. The \emph{deletion metric} of Petsiuk et
al.~\cite{petsiuk2018rise} removes pixels in decreasing order of attributed
importance and records how quickly the model's confidence falls, comparing
against removal in random order. A faithful map degrades the prediction faster
than chance.

The metrics used to compare two explanations with each other are borrowed from
the saliency-model literature, where they were developed to compare a predicted
map against human eye fixations~\cite{bylinskii2019metrics}: Pearson
correlation, structural similarity~\cite{wang2004ssim}, and the
intersection-over-union of the highest-scoring regions. Bylinskii et al. observe
that rankings produced by different metrics frequently disagree, which is the
argument for reporting several rather than choosing one.

A separate line of work asks whether attribution methods are testing anything at
all. Adebayo et al.~\cite{adebayo2018sanity} randomize model weights
progressively and verify that explanations change in response; a map that
survives the destruction of the model it purports to explain is not explaining
it. The perturbation experiment of Section~\ref{sec:perturbation} is a graded
form of that test, run in the opposite direction.

\section{Gravitational-wave glitch classification}

The LIGO interferometers are limited by non-Gaussian transient noise of
instrumental and environmental origin. These glitches occur at rates of several
per minute, can mimic astrophysical signals, and are grouped into morphological
families that correlate with distinct physical causes.

The Gravity Spy project~\cite{zevin2017, bahaadini2018} combines machine
classification with citizen science to label them. The release used
here~\cite{glanzer2021} covers observing runs O1 through O3b and provides, for
each Omicron trigger, the trigger metadata, the per-class confidence of a
convolutional classifier, the assigned label, and Omega-scan spectrograms.

The classification task is therefore image classification over time--frequency
representations. Its interest for this thesis is not that it is difficult ---
with four well-separated classes it is not --- but that the explanations have a
consumer. Glitch classification feeds detector characterization, in which the
morphology of a glitch family is traced back to an instrumental cause and the
instrument is modified. An explanation that points at the wrong part of the
spectrogram sends a commissioner to the wrong subsystem.

\section{The gap this thesis occupies}

The two closest works can now be placed precisely. Rabbi et
al.~\cite{rabbi2026} measure explanation fidelity under quantization, but for
conventional networks, post-training weight-only quantization, and purely
spatial explanations. Smith et al.~\cite{smith2026} address spiking networks and
separate the two quantization targets, but measure divergence in \emph{population
firing statistics} aggregated over a dataset, not in explanations.

The unoccupied ground is per-sample, spatially and temporally resolved
explanation fidelity for a spiking classifier with the two quantization targets
separated. That is the study this thesis specifies.

Neither work, however, establishes what its metrics report when nothing has
changed. Rabbi et al. present similarity values around $0.85$ as evidence of
preserved explanations; Smith et al. present EMD values between $0.001$ and
$0.15$ as evidence of behavioural divergence. Chapter~\ref{ch:results} shows
that for the architecture studied here the corresponding floors are $0.975$ for
numerical noise and $0.153$ for models differing only in initialization, so a
value of $0.85$ falls between the two and cannot be interpreted without knowing
which comparison produced it. Establishing those floors is the contribution with
the widest applicability.
